\documentclass{article}

\usepackage{float}
\restylefloat{table}

\usepackage{booktabs}

\title{Team Contributions: Final\\\progname}

\author{\authname}

\date{}

\input{../Comments}
\input{../Common}

\begin{document}

\maketitle

This document summarizes the contributions of each team member for the final
demonstration and documentation.  The time period of interest is the time
between Rev 0 and the Final documentation; the contributions prior to Rev0 are
NOT included.

\section{Team Meeting Attendance}

\begin{table}[H]
\centering
\begin{tabular}{ll}
\toprule
\textbf{Student} & \textbf{Meetings}\\
\midrule
Total & 7\\
Mahad Ahmed & 7\\
Abyan Jaigirdar & 7\\
Prerna Prabhu & 7\\
Farhan Rahman & 7\\
Ali Zia & 7\\
\bottomrule
\end{tabular}
\end{table}

\noindent
The team arranges meetings almost weekly, to clarify objectives and align on priorities on major deliverables, with additional meetings scheduled as needed, typically in the days leading up to major deadlines. We have held a total of seven meetings in the time period of interest. Every team member has been able to attend these meetings, even if occasionally arriving slightly late or needing to leave early.

\section{Supervisor/Stakeholder Meeting Attendance}

\noindent \textbf{Supervisor's Name: } Luke Schuurman

\begin{table}[H]
\centering
\begin{tabular}{ll}
\toprule
\textbf{Student} & \textbf{Meetings}\\
\midrule
Total & 1\\
Abyan & 1\\
Ali & 1\\
Farhan & 1\\
Mahad & 1\\
Prerna & 1\\
\bottomrule
\end{tabular}
\end{table}

\noindent
Supervisor meetings were arranged when objectives and priorities needed to be clarified and all members attended those meetings. This period of interest we diverged from the original capstone plan and only decided to have one meeting to discuss the action plan and approach the team was going to take. 

\section{Lecture Attendance}

Between the Rev 0 and the Final demos, there has been approximately 1 lecture held as part of the course schedule.


\begin{table}[H]
\centering
\begin{tabular}{ll}
\toprule
\textbf{Student} & \textbf{Lectures}\\
\midrule
Total & 1\\
Mahad Ahmed & 0\\
Abyan Jaigirdar & 0\\
Prerna Prabhu & 1\\
Farhan Rahman & 0\\
Ali Zia & 0\\
\bottomrule
\end{tabular}
\end{table}

Due to timing of the lecture and scheduling conflicts, only one team member who lives closest to the university was able to attend the lecture. The team member who attended the lecture shared the information with the rest of the team and we had a discussion about it in our next team meeting.

\section{TA Document Discussion Attendance}


\noindent \textbf{TA's Name: } Tiago de Moraes Machado

\begin{table}[H]
\centering
\begin{tabular}{ll}
\toprule
\textbf{Student} & \textbf{Lectures}\\
\midrule
Total & 1\\
Mahad Ahmed & 0\\
Abyan Jaigirdar & 0\\
Prerna Prabhu & 0\\
Farhan Rahman & 0\\
Ali Zia & 0\\
\bottomrule
\end{tabular}
\end{table}

\noindent During the time period between the Rev 0 and Final document deadlines, there was one TA document discussion session. None of the members were able to attend the TA document discussion session due to scheduling conflicts, religious reasons, and injuries.

\section{Commits}

\begin{table}[H]
\centering
\begin{tabular}{lll}
\toprule
\textbf{Student} & \textbf{Commits} & \textbf{Percent}\\
\midrule
Total & 201 & 100\% \\
Mahad Ahmed & 70 & 20.40\% \\
Abyan Jaigirdar & 40 & 16.92\% \\
Prerna Prabhu & 87 & 23.87\% \\
Farhan Rahman & 88 & 18.91\% \\
Ali Zia & 42 & 19.90\% \\
\bottomrule
\end{tabular}
\end{table}

\section{Issue Tracker}

\begin{table}[H]
\centering
\begin{tabular}{ll}
\toprule
\textbf{Student}  & \textbf{Issues Assigned}\\
\midrule
Mahad Ahmed  & 23 \\
Abyan Jaigirdar  & 15 \\
Prerna Prabhu  & 28 \\
Farhan Rahman  & 16 \\
Ali Zia & 16 \\
\bottomrule
\end{tabular}
\end{table}

\noindent
One team member created and opened all issues for upcoming milestones and document sections. This was done intentionally to maintain a consistent issue structure and keep everything organized. During team meetings, these pre-created issues were reviewed based on how much time it would take and then assigned to team members evenly. \\

\noindent
Because issue creation was done in this way, the column "O+C" (Opened and Closed Issues) does not accurately reflect the individual contribution. Instead, the key metric is just the number of issues assigned to each team member. Therefore, we have chosen to remove the "O+C" column and report only the Assigned (Closed) issues.

\section{Team Charter Trigger Items}

\subsection{Summary of Triggers}
The team charter established several quantified triggers to help maintain accountability and ensure consistent contribution throughout the project. These include:
\begin{itemize}
    \item \textbf{Attendance:} Members are expected to maintain a 100\% attendance rate for all scheduled meetings unless an acceptable excuse (such as a health issue or family emergency) is provided.
    \item \textbf{GitHub Activity:} Each member must demonstrate consistent GitHub activity every week, with issues actively in progress or commits made for review.
    \item \textbf{Task Completion:} All assigned work must be completed and delivered on time according to deadlines set during weekly meetings.
\end{itemize}

\subsection{Trigger Violations}
Throughout the period between Rev 0 and the Final demonstration, the team has not experienced any major violations of the triggers outlined in the charter. All members maintained full attendance at team meetings, continued to show consistent GitHub activity, and met their assigned deadlines. As in the previous period, minor delays in individual tasks occurred on occasion, but these were communicated promptly and resolved collaboratively without affecting the overall project timeline or deliverables.

\subsection{Plan to Address Violations}
As the project is now concluding, the three-step escalation process defined in the charter was not needed at any point during the project:
\begin{enumerate}
    \item First incident: verbal reminder during team meeting.
    \item Second incident: discussion with the TA to address underlying issues.
    \item Third or repeated incidents: escalation to the course instructor.
\end{enumerate}

Overall, the triggers proved to be appropriately calibrated for the team. The expectations around attendance, GitHub activity, and task completion were realistic and effective in maintaining accountability without being overly restrictive. The team found that open communication and early reporting of potential delays were sufficient to prevent any formal violations from occurring throughout the entire duration of the project.

\section{Additional Productivity Metrics}

\wss{If your team has additional metrics of productivity, please feel free to
add them to this report.}

\end{document}